\scanpage{24}
\deleted{If $x\in X\setminus Y$, $\exists\alpha$ such that $f_\alpha(x)\ne0$, and}

define $Y$ as the set of common zeros. \deleted{of $X$} : these functions, moreover they separate the points of $X\setminus Y$, for if $x,y$ are distinct points, \deleted{$\exists\beta$ such that $f_\alpha(x)\ne0$ $g_\beta(x)\ne g_\beta(y)$ whence}

If the $f_\alpha$ do \unsure{not} suffice to separate them, let us show that the $f_\alpha g_\beta$ separate them : indeed, $\exists\beta$ such that $g_\beta(x)\ne g_\beta(y)$ \deleted{whence} and $\alpha$ such that $\lambda=f_\alpha(x)\ne0$, thus $f_\alpha(y)$ $(=f_\alpha(x)=\lambda\text{ by hypothesis})\ne0$, whence
\[
f_\alpha\underbrace{g_\beta(x)}_{\lambda}=\lambda g_\beta(x)\ne f_\alpha g_\beta(y)=\lambda g_\beta(y)
\]

Thus :

\uline{Proposition}. If we have M1 : M\unsure{4|7} and M11, then for every \deleted{\unsure{such|closed}} \inserted{tame} subset $Y$ of a tame space $X$, which is the set of common zeros of a finite family of tame functions on $X$, one can find such a family that separates the points of $X$ (thus $X/Y$ is a tame space).

\mspage{23}
